\section*{Topic 2 Completion Audit and Validation Protocol}
\phantomsection
\addcontentsline{toc}{section}{Topic 2 Completion Audit and Validation Protocol}

This supplement distinguishes corrected analysis from experiments that remain
prospective. The proposal asks for robust multilingual meeting transcription,
adaptation and a measured accuracy--latency trade-off. The existing 468 canonical
transcriptions support a component-level pilot; they do not establish all three
requirements in real meetings. No additional ASR experiment is claimed in this
supplement.

\subsection*{Audit of the Extended Live Log}
The historical 22-observation study must be distinguished from the current
\texttt{results/\allowbreak live\_\allowbreak latency.csv}, which contains 38 observations. Inferring
session boundaries from the utterance counter restarting at one yields three
sessions. Excluding those three first observations gives 35 observations with
median 0.575\,s, p95 1.614\,s and maximum 11.021\,s. The third inferred session
alone contains 15 non-initial observations, with median 0.657\,s and p95
4.814\,s. The slow observation is retained. The legacy log does not record
session timestamps, model settings, hardware or concurrent load, so pooling
these observations is descriptive and does not establish comparable conditions.
The first observation is a cold-start proxy rather than an isolated measurement
of model loading. The audit is reproducible with
\texttt{python -m backend.\allowbreak evaluation.\allowbreak topic2\_\allowbreak readiness}; its report records the
input file's SHA-256 digest.

Code inspection further bounds the measurement: the elapsed time is recorded
before WebSocket transmission, rather than at browser receipt or rendering.
Empty or rejected outputs can return before logging. Consequently the CSV
neither measures end-to-end display latency nor estimates the overall failure
rate. The historical p95 of 0.844\,s is not a service guarantee. Adding the
configured silence interval is only an approximation of part of the additional
delay.

\subsection*{Meeting and Speaker Validation}
The simulated meeting builder alternates speaker labels across independent
telephone clips without verified persistent speaker identities. Its DER cannot
validate real speaker-attribution accuracy. A complementary experiment requires
licensed meeting audio with verified transcripts, speaker identities and overlap
annotations. Development and evaluation must be separated by meeting, with
channel choice, time windows and scoring rules fixed before inference. WER,
DER and an established speaker-attributed transcription metric should be
reported per meeting, including failures. Uncertainty should respect meeting
boundaries. An English-only extension would not validate multilingual meetings.
These data and results are not yet supplied by the existing simulated corpus.

\subsection*{Domain and Multilingual Validation}
A vocabulary-specific evaluation should compare identical recordings with and
without a fixed glossary prompt, holding the model and decoding settings
constant. Human-verified references should include target terms and matched
negative controls, allowing both term recall and hallucinated terms to be
reported alongside WER. Any additional adaptation must use separate training
material. The existing banking-domain LoRA result cannot establish a benefit
for Nestl\'e terminology. Likewise, separate English, French and German clips
do not establish within-utterance language-switching performance. An annotated
switching set and additional documented acoustic conditions are needed; their
results must remain separate from the canonical 8\,dB white-noise campaign.

\subsection*{Controlled Live Validation}
A prospective live study should replay held-out audio at real-time speed through
the WebSocket path, varying model size and endpointing settings on identical
content with repeated, randomised trials. Each run must record configuration,
hardware, reference and final hypothesis, speech end, endpoint decision,
server transmission and browser receipt/render durations. Monotonic durations
must not be formed by subtracting unsynchronised client and server clocks.
Empty outputs, rejected outputs, failures and cold starts must remain visible.
Joint WER, median/p95/maximum latency and failure rates are required before
selecting an accuracy--latency operating point. A simulated Teams trial must
also document its actual audio-routing path; an adjacent browser prototype
alone is not a Teams integration. These protocols define remaining validation
work, not completed experiments or measured accessibility benefits.
